\paragraph{線形代数の練習問題}行列$\bm{A}=\begin{pmatrix}1 & 3 \\2 & 4\end{pmatrix},\bm{B}=\begin{pmatrix}3&1&5\\2&4&6\end{pmatrix},\bm{C}=\begin{pmatrix}4&1\\2&3\end{pmatrix}$，列ベクトル$\bm{x}=\begin{pmatrix}1\\2\\3\end{pmatrix}$,行ベクトル$\bm{y}=\begin{pmatrix}2&8\end{pmatrix}$とするとき，次の計算をRで実行するコードを書きなさい。なお，計算できないものについては「計算できない」としてコードは書かないこと。
\begin{enumerate}
	\item $\bm{A}+\bm{B}$
	\item $\bm{A}-\bm{C}$
	\item $\bm{AB}$
	\item $\bm{AC}$
	\item $\bm{B}'\bm{A}$
	\item $\bm{A}\bm{y}'$
	\item $\bm{xy}$
	\item $\bm{xB}$
	\item $\bm{x}'\bm{B}'$
	\item $\bm{yx}$
\end{enumerate}

\paragraph{連立方程式を解く}次の連立方程式をRで解きなさい。
\[
	\left\{\begin{array}{rr}x-2 y+3 z= & 1 \\ 3 x+y-5 z= & -4 \\ -2 x+6 y-9 z= & -2\end{array}\right.
\]

\end{document}
